\begin{name}
	{Biên soạn: Thầy Nguyễn Hoài Nam \& Phản biện: Thầy Mui Doan}
	{Đề chuẩn hóa HK2 trường THPT Nguyễn Khuyến, năm 2022 - 2023}
\end{name}
\cautn
\Opensolutionfile{ans}[ans/ans-cau-TN-1]
\begin{ex}%[De-chuan-hoa-so-3]%[Nguyễn Hoài Nam]%[0D2N1-2]
	Bất phương trình $-x^2+2x+3<0$ có tập nghiệm là
	\choice
	{$(-1;3)$}
	{$(-3;1)$}
	{$[-1;3]$}
	{\True$(-\infty;-1) \cup (3;+\infty)$}
	\loigiai{
		Ta có $-x^2+2x+3<0\Leftrightarrow \hoac{&x<-1\\&x>3.}$ \\
		Vậy tập nghiệm của bất phương trình là $S=(-\infty;-1)\cup (3;+\infty)$.
	}
\end{ex}
\begin{ex}%[De-chuan-hoa-so-3]%[Nguyễn Hoài Nam]%[0H9H3-3]
	Xét vị trí tương đối của hai đường thẳng $d_1\colon x-2y+1=0$ và $d_2\colon -3x+6y-10=0$
	\choice
	{\True Song song}
	{Cắt nhau nhưng không vuông góc}
	{Vuông góc với nhau}
	{Trùng nhau}
	\loigiai{
		Ta có $\dfrac{1}{-3}=\dfrac{-2}{6}\ne\dfrac{1}{10}$ nên $d_1$ song song với $d_2$.
	}
\end{ex}
\begin{ex}%[De-chuan-hoa-so-3]%[Nguyễn Hoài Nam]%[0H9N4-1]
	Cho đường tròn $\left(C\right)\colon (x-2)^2+(y+6)^2=8$ có tâm $I$ và bán kính là
	\choice
	{\True $I(2;-6),R=2\sqrt{2}$}
	{$I(2;-6),R=2$}
	{$I(2;6),R=2\sqrt{2}$}
	{$I(2;-6),R=8$}
	\loigiai{
		$(C)$ có tâm $I(2;-6)$ và bán kính $R=\sqrt{8}=2\sqrt{2}$
	}
\end{ex}

\begin{ex}%[De-chuan-hoa-so-3]%[Nguyễn Hoài Nam]%[0H9N3-2]
	Cho đường thẳng $d\colon-3x+2y-4=0$. Một véc-tơ pháp tuyến của đường thẳng là
	\choice
	{$\overrightarrow{n}=\left(3;2\right)$}
	{$\overrightarrow{n}=\left(-3;-2\right)$}
	{\True $\overrightarrow{n}=\left(-3;2\right)$}
	{$\overrightarrow{n}=\left(2;3\right)$}
	\loigiai{Đường thẳng $d$ nhận $\overrightarrow{n}=(-3;2)$ làm một véc-tơ pháp tuyến.
	}
\end{ex}
\begin{ex}%[De-chuan-hoa-so-3]%[Nguyễn Hoài Nam]%[0D2N1-2]
	Bất phương trình $2x^2-5x+2\geq 0$ có tập nghiệm là
	\choice
	{$\left(-\infty;\dfrac{1}{2} \right]$}
	{$\left(-\infty;\dfrac{1}{2}\right)\cup \left(2;+\infty \right)$}
	{\True $\left(-\infty;\dfrac{1}{2}\right]\cup \left[2;+\infty    \right)$}
	{$\left[\dfrac{1}{2};2\right]$}
	\loigiai{
		Ta có $2x^2-5x+2\geq 0\Leftrightarrow\hoac{&x\geq2\\&x\leq \dfrac{1}{2}.}$\\ Vậy tập nghiệm của phương trình là $S=\left(-\infty;\dfrac{1}{2}\right]\cup \left[2;+\infty \right)$.
	}
\end{ex}
\begin{ex}%[De-chuan-hoa-so-3]%[Nguyễn Hoài Nam]%[0D0H1-2]
	Gieo ngẫu nhiên một đồng tiền cân đối, đồng chất $5$ lần. Tính số phần tử không gian mẫu.
	\choice
	{$10$}
	{\True $32$}
	{$16$}
	{$64$}
	\loigiai{
		Ta có $n\left(\Omega\right)=2^5=32$.
	}
\end{ex}
\begin{ex}%[De-chuan-hoa-so-3]%[Nguyễn Hoài Nam]%[0H9H3-5]
	Khoảng cách từ điểm $A(1;1)$ đến đường thẳng $5x-12y-6=0$ là
	\choice
	{$-13$}
	{\True $1$}
	{$-1$}
	{$64$}
	\loigiai{
		Đặt $\Delta\colon  5x-12y-6=0$.\\
		Khi đó $\mathrm{d}(A,\Delta)=\dfrac{\left|5\cdot1-12\cdot1-6\right|}{\sqrt{5^2+\left(-12\right)^2}}=1$.
	}
\end{ex}
\begin{ex}%[De-chuan-hoa-so-3]%[Nguyễn Hoài Nam]%[0D0H2-2]
	Gieo con súc sắc hai lần. Gọi $A$ là biến cố để sau hai lần gieo có ít nhất một mặt $6$ chấm
	\choice
	{$A=\left\{\left(6;1\right),\left(6;2\right),\left(6;3\right),\left(6;4\right),\left(6;5\right)\right\}$}
	{$A=\left\{\left(1;6\right),\left(2;6\right),\left(3;6\right),\left(4;6\right),\left(5;6\right),\left(6;6\right)\right\}$}
	{\True $A=\left\{\left(1;6\right),\left(2;6\right),\left(3;6\right),\left(4;6\right),\left(5;6\right),\left(6;6\right),\left(6;1\right),\left(6;2\right),\left(6;3\right),\left(6;4\right),\left(6;5\right)\right\}$}
	{$A=\left\{\left(1;6\right),\left(2;6\right),\left(3;6\right),\left(4;6\right),\left(5;6\right)\right\}$}
	\loigiai{
		Để sau 2 lần có ít nhất 1 mặt $6$ chấm ta có $2$ trường hợp sau:\\
		\textbf{Trường hợp 1:} Lần $1$ xuất hiện mặt $6$ chấm, lần $2$ tùy ý.\\
		\textbf{Trường hợp 2:} Lần $2$ xuất hiện mặt $6$ chấm, lần $1$ tùy ý.
	}
\end{ex}
\begin{ex}%[De-chuan-hoa-so-3]%[Nguyễn Hoài Nam]%[0H9H3-4]
	Góc giữa hai đường thẳng $\Delta_{1}\colon3x-2y+1=0$;  $\Delta_{2}\colon\heva{&x=t\\&y=7-5t}$ $(t\in \mathbb{R})$.
	\choice
	{\True $45^\circ $}
	{$30^\circ $}
	{$135^\circ $}
	{$60^\circ $}
	\loigiai{
		$\Delta_{1}$ có $1$ véc-tơ pháp tuyến là $\overrightarrow{n}_{1}=\left(3;-2\right)$.\\
		$\Delta_{2}$ có $1$ véc-tơ chỉ phương là $\overrightarrow{u}=\left(1;-5\right)$ nên có $1$ véc-tơ pháp tuyến là $\overrightarrow{n}_{2}=\left(5;1\right)$.\\
		Do đó $\cos(\Delta_{1},\Delta_{2})=\dfrac{\left|\overrightarrow{n}_{1}\cdot\overrightarrow{n}_{2}\right|}{\left|\overrightarrow{n}_{1}\right|\cdot\left|\overrightarrow{n}_{2}\right|}=\dfrac{\sqrt{2}}{2} \Rightarrow \left(\Delta_{1},\Delta_{2}\right)=45^\circ$.
	}
\end{ex}
\begin{ex}%[De-chuan-hoa-so-3]%[Nguyễn Hoài Nam]%[0D8H2-5]
	Trên bàn có $8$ cây bút chì khác nhau, $6$ cây bút bi khác nhau và $10$ cuốn tập khác nhau. Số cách khác nhau để chọn được đồng thời một cây bút chì, một cây bút bi và một cuốn tập.
	\choice
	{$48$}
	{$60$}
	{$24$}
	{\True $480$}
	\loigiai{
		\textbf{Công đoạn $1$:} Chọn $1$ cây bút chì từ $8$ cây bút chì có $\mathrm{C}_{8}^1=8$ cách.\\
		\textbf{Công đoạn $2$:} Chọn $1$ cây bút bi từ $6$ cây bút chì có $\mathrm{C}_{6}^1=6$ cách.\\
		\textbf{Công đoạn $3$:} Chọn $1$ cây cuốn tập từ $10$ cuốn tập có $\mathrm{C}_{10}^1=10$ cách.\\
		Theo quy tắc nhân có $8\cdot6\cdot10=480$ cách.
	}
\end{ex}
\begin{ex}%[De-chuan-hoa-so-3]%[Nguyễn Hoài Nam]%[0H9H5-1]
	Elip $(E)\colon\dfrac{x^2}{25}+\dfrac{y^2}{16}=1$ có tiêu cự bằng
	\choice
	{\True $6$}
	{$3$}
	{$9$}
	{$18$}
	\loigiai{
		Gọi $(E)$ có dạng $\dfrac{x^2}{a^2}+\dfrac{y^2}{b^2}=1,\left(a>b>0\right)$. Và tiêu cự $F_{1}F_{2}=2c, \left(c>0\right)$. \\Khi đó $c^2=a^2-b^2=9\Rightarrow c=3$.\\
		Vậy $F_{1}F_{2}=2c=6$.
	}
\end{ex}
\begin{ex}%[De-chuan-hoa-so-3]%[Nguyễn Hoài Nam]%[0D8H3-2]
	Viết khai triển theo công thức nhị thức Niu-tơn $\left(x-y\right)^5$ là
	\choice
	{$x^5+5x^4y+10x^3y^2+10x^2y^3+5xy^4+y^5$}
	{$x^5-5x^4y-10x^3y^2-10x^2y^3-5xy^4+y^5$}
	{\True $x^5-5x^4y+10x^3y^2-10x^2y^3+5xy^4-y^5$}
	{$x^5+5x^4y-10x^3y^2+10x^2y^3-5xy^4+y^5$}
	\loigiai{
		Ta có
		\begin{eqnarray*}
			&\left(x-y\right)^5&=\mathrm{C}_{5}^0x^5\left(-y\right)^0+\mathrm{C}_{5}^1x^4\left(-y\right)^1+\mathrm{C}_{5}^2x^3\left(-y\right)^2+\mathrm{C}_{5}^3x^2\left(-y\right)^3+\mathrm{C}_{5}^4x^1\left(-y\right)^4+\mathrm{C}_{5}^5x^0\left(-y\right)^5\\
			&&=x^5-5x^4y+10x^3y^2-10x^2y^3+5xy^4-y^5.
		\end{eqnarray*} 
	}
\end{ex}


\Closesolutionfile{ans}
\indapan{12}{ans/ans-cau-TN-1}

\cauds
\Opensolutionfile{ans}[ans/ans-cau-DS-1]
\begin{ex}%[De-chuan-hoa-so-3]%[Nguyễn Hoài Nam]%[0D8H3-4]
	Cho nhị thức $(x+2)^5$, xét tính đúng sai của các mệnh đề sau.
	\choiceTF[t]
	{\True Khai triển $(x+2)^5$ có $6$ số hạng}
	{Hệ số của số hạng chứa $x^2$ là $40$}
	{\True Tổng hệ số trong khai triển là $243
		$ }
	{Hệ số của số hạng không chứa $x$ là $1$}
	\loigiai{
		\begin{itemchoice}
			\itemch Đúng. Trong khai triển $(a+b)^n$ có $n+1$ số hạng. Vậy khai triển $\left(x+2\right)^5$ có 6 số hạng.
			\itemch Sai. Vì
			$$(x+2)^5=\mathrm{C}_5^0x^5+\mathrm{C}_5^1x^4\cdot 2^1+\mathrm{C}_5^2\cdot x^3\cdot 2^2+\mathrm{C}_5^3 x^2\cdot 2^3+\mathrm{C}_5^4 x^4\cdot 2^4+\mathrm{C}_5^5 2^5=x^5+10x^4+40x^3+80x^2+80x+32$$
			\itemch Đúng. Ta có 
			$$a_0+a_1+a_2+a_3+a_4+a_5=\mathrm{C}_5^0+\mathrm{C}_5^1\cdot 2^1+\mathrm{C}_5^2\cdot 2^2+\mathrm{C}_5^3 \cdot 2^3+\mathrm{C}_5^4\cdot 2^4+\mathrm{C}_5^5=243$$
			\itemch Sai. Hệ số của số hạng không chứa $x$ là $\mathrm{C}_5^5\cdot 2^5$. 
		\end{itemchoice}
	}
\end{ex}
\begin{ex}%[De-chuan-hoa-so-3]%[Nguyễn Hoài Nam]%[0H9H3-2]
	Trong mặt phẳng với hệ trục tọa độ $Oxy$, cho hai điểm $A(3;-2)$ và $B(1,2)$. Gọi $d$ là đường trung trực của đoạn thẳng $AB$.
	\choiceTF[t]
	{\True Tọa độ trung điểm $I$ của đoạn thẳng $AB$ là $(2;0)$}
	{Một véc-tơ pháp tuyến của đường thẳng $d$ có tọa độ là $(2;4)$}
	{\True Một véc-tơ chỉ phương của đường thẳng $d$ có tọa độ là $(2;1)$}
	{Phương trình tham số của $d$ là $\heva{&x=2+2t \\ &y=1}$ ($t\in\mathbb{R}$)}
	\loigiai{
		\begin{itemchoice}
			\itemch Đúng. Trung điểm $I$ của đoạn thẳng $AB$ có tọa độ là $\left( \dfrac{3+1}{2}; \dfrac{-2+2}{2} \right) = (2;0)$.
			\itemch Sai. Ta có $\overrightarrow{AB}=(-2;4)$ là một véc-tơ pháp tuyến của đường thẳng $d$. Véc-tơ có tọa độ $(2;4)$ không cùng phương với $\overrightarrow{AB}$ nên không phải là một véc-tơ pháp tuyến của đường thẳng $d$.
			\itemch Đúng. Đường thẳng $d$ có một véc-tơ pháp tuyến là $\overrightarrow{AB} = (-2;4)$, suy ra có một véc-tơ chỉ phương là $\overrightarrow{u} = (2;1)$.
			\itemch Sai. Do đường thẳng $d$ có véc-tơ chỉ phương là $\overrightarrow{u}=(2;1)$ và đi qua điểm $I(2;0)$, nên $d$ có phương trình tham số là $\heva{&x=2+2t \\ &y=t}$ ($t\in\mathbb{R}$).
		\end{itemchoice}
	}	
\end{ex}
\begin{ex}%[De-chuan-hoa-so-3]%[Nguyễn Hoài Nam]%[0D0H2-4]
	Một nhóm gồm $6$ học sinh nữ và $4$ học sinh nam. 
	\choiceTF[t]
	{Chọn ngẫu nhiên $3$ học sinh có $10$ cách}
	{Chọn $2$ học sinh nam và $1$ học sinh nữ có $12$ cách}
	{\True Chọn ngẫu nhiên $3$ học sinh, xác suất để chọn được $2$ học sinh nữ và $1$ học sinh nam bằng $\dfrac{1}{2}$}
	{\True Chọn ngẫu nhiên $3$ học sinh, xác suất để $3$ học sinh được chọn đều là nam bằng $\dfrac{1}{30}$ }
	\loigiai{
		\begin{itemchoice}
			\itemch Sai. Số cách chọn ngẫu nhiên $3$ học sinh là $\mathrm{C}_{10}^3=120$ cách.
			\itemch Sai. Số cách chọn $2$ nam và $1$ nữ là $\mathrm{C}_4^2 \cdot \mathrm{C}_6^1=36$ cách.  
			\itemch Đúng. Số phần tử không gian mẫu $n\left(\Omega\right)= \mathrm{C}_{10}^3=120$.\\
			Gọi $A$ là biến cố \lq\lq Chọn được $3$ học sinh gồm $2$ học sinh nữ và $1$ học sinh nam\rq\rq.\\
			Số kết quả thuận lợi của biến cố $A$ là $n(A)=\mathrm{C}_{6}^2\cdot\mathrm{C}_{4}^1=60$.\\
			Vậy xác suất cần tìm là $P(A)=\dfrac{n(A)}{n(\Omega)}=\dfrac{1}{2}$.
			\itemch Đúng. Số phần tử không gian mẫu: $n(\Omega)= \mathrm{C}_{10}^3=120$.\\
			Gọi $A$ là biến cố \lq\lq $3$ học sinh được chọn đều là nam \rq\rq.\\
			Số kết quả thuận lợi của biến cố $A$ là $n(A)=\mathrm{C}_{4}^3=4$.\\
			Vậy xác suất cần tìm là $P(A)=\dfrac{n(A)}{n(\Omega)}=\dfrac{1}{30}$.
		\end{itemchoice}
	}	
\end{ex}
\begin{ex}%[De-chuan-hoa-so-3]%[Nguyễn Hoài Nam]%[0H9H5-1]
	Cho elip $(E)$ có độ dài trục lớn bằng $20$ và tiêu cự bằng $10$.
	\choiceTF[t]{\True Nếu $M$ là một điểm tùy ý trên elip $(E)$ thì $MF_1+MF_2=20$ với $F_1$, $F_2$ là hai tiêu điểm của elip $(E)$}
	{\True Elip $(E)$ có hai tiêu điểm là $F_1(-5;0)$ và $F_2(5;0)$}
	{Elip $(E)$ có độ dài trục bé là $150$.}
	{\True Phương trình chính tắc của elip $(E)$ là $\dfrac{x^2}{100}+\dfrac{y^2}{75}=1$}
	\loigiai{
		\begin{itemchoice}
			\itemch Đúng. Do độ dài trục lớn là $2a=20$ nên theo định nghĩa ta có $MF_1+MF_2=2a=20$.
			\itemch Đúng. Do tiêu cự $2c=10$ nên $c=5$. Suy ra, elip $(E)$ có hai tiêu điểm là $F_1(-5;0)$ và $F_2(5;0)$.
			\itemch Sai. Ta có $c^2=a^2-b^2\Leftrightarrow b^2=a^2-c^2=10^2-5^2=75$. Suy ra, độ dài trục bé bằng $2b=10\sqrt{3}$.
			\itemch Đúng. Phương trình chính tắc của elip cần tìm là
			$$ (E)\colon \dfrac{x^2}{100}+\dfrac{y^2}{75}=1.$$
		\end{itemchoice}
	}
\end{ex}

\Closesolutionfile{ans}
\indapan{4}{ans/ans-cau-DS-1}

\caukq
\Opensolutionfile{ans}[ans/ans-cau-KQ-1]


\begin{ex}%[De-chuan-hoa-so-3]%[Nguyễn Hoài Nam]%[0D8V2-4]
	Ở một Đoàn trường phổ thông có $5$ thầy giáo, $4$ cô giáo và $8$ học sinh. Có $X$ cách chọn ra một đoàn công tác gồm $7$ người trong đó có $1$ trưởng đoàn là thầy giáo, $1$ phó đoàn là cô giáo và đoàn công tác có ít nhất $4$ học sinh. Tính tổng các chữ số của $X$.
	\shortans[oly]{$12$}	
	\loigiai{
		\begin{itemize}
			\item Trường hợp $1$: Đoàn công tác gồm $1$ thầy giáo, $1$ cô giáo và $5$ học sinh.
			\begin{itemize}
				\item Chọn ra $1$ thầy giáo từ $5$ thầy giáo có $\mathrm{C}_5^1$ cách.
				\item Chọn ra $1$ cô giáo từ $4$ cô giáo có $\mathrm{C}_4^1$ cách.
				\item Chọn ra $5$ học sinh từ $8$ học sinh có $\mathrm{C}_8^5$ cách.
			\end{itemize}
			Do chỉ có $1$ thầy giáo và $1$ cô giáo nên chỉ có $1$ cách chọn trưởng đoàn và phó đoàn.\\
			Theo quy tắc nhân có $\mathrm{C}_5^1\cdot\mathrm{C}_4^1\cdot\mathrm{C}_8^5\cdot1=1 \, 120$ cách.
			\item Trường hợp $2$: Đoàn công tác gồm $2$ thầy giáo, $1$ cô giáo và $4$ học sinh.
			\begin{itemize}
				\item Chọn ra $2$ thầy giáo từ $5$ thầy giáo có $\mathrm{C}_5^2$ cách.
				\item Chọn ra $1$ cô giáo từ $4$ cô giáo có $\mathrm{C}_4^1$ cách.
				\item Chọn ra $4$ học sinh từ $8$ học sinh có $\mathrm{C}_8^4$ cách.
				\item Chọn $1$ thầy giáo từ $2$ thầy giáo làm trưởng đoàn có $\mathrm{C}_2^1$ cách.
			\end{itemize}
			Theo quy tắc nhân có $\mathrm{C}_5^2\cdot\mathrm{C}_4^1\cdot\mathrm{C}_8^4\cdot\mathrm{C}_2^1=5\,600$ cách.
			\item Trường hợp $3$: Đoàn công tác gồm $1$ thầy giáo, $2$ cô giáo và $4$ học sinh.
			\begin{itemize}
				\item Chọn ra $1$ thầy giáo từ $5$ thầy giáo có $\mathrm{C}_5^1$ cách.
				\item Chọn ra $2$ cô giáo từ $4$ cô giáo có $\mathrm{C}_4^2$ cách.
				\item Chọn ra $4$ học sinh từ $8$ học sinh có $\mathrm{C}_8^4$ cách.
				\item Chọn $1$ cô giáo từ $2$ cô giáo làm phó đoàn có $\mathrm{C}_2^1$.
			\end{itemize}
			Theo quy tắc nhân có $\mathrm{C}_5^1\cdot\mathrm{C}_4^2\cdot\mathrm{C}_8^4\cdot\mathrm{C}_2^1=4\, 200$ cách.			
		\end{itemize}		
	Theo quy tắc cộng ta có $1\,120+5\,600+4\,200=10\,920$ cách.
	}
\end{ex}
\begin{ex}%[De-chuan-hoa-so-3]%[Nguyễn Hoài Nam]%[0D0V2-3]
	Mật khẩu mở máy tính của Dũng gồm $8$ kí tự đôi một khác nhau, trong đó có $2$ kí tự đầu là chữ số, $6$ kí tự sau là các chữ cái thuộc tập hợp $\left\{ A,B,C,D,E,F \right\}$. Không may Dũng quên mất $3$ kí tự đầu tiên. Dũng chọn ra $2$ chữ số và một chữ cái thuộc tập hợp trên một cách ngẫu nhiên và thử mở máy tính. Xác suất để Dũng mở được máy tính có dạng $\dfrac{a}{b}$. Tính $a+b$.
	\shortans[oly]{$541$}
	
	\loigiai{
		\textbf{Công đoạn $1$:} Chọn $2$ số từ $10$ số và có phân biệt thứ tự nên có $\mathrm{A}_{10}^2$ cách.\\
		\textbf{Công đoạn $2$:} Chọn ra $1$ chữ từ $6$ chữ cái có $\mathrm{C}_6^1$ cách.\\
		Vậy theo quy tắc nhân có cách lập mật khẩu là $\mathrm{A}_{10}^2\cdot\mathrm{C}_6^1=540$ cách.\\
		Vậy số phần tử không gian mẫu là $ n(\Omega)=540$.\\
		Gọi biến cố $A$ \lq\lq Xác xuất để Dũng mở được máy\rq\rq.\\
		Do chỉ có $1$ mật mã đúng duy nhất nên $n(A)=1$.\\
		Vậy $P(A)=\dfrac{n(A)}{n(\Omega)}=\dfrac{1}{540}$.
	}
\end{ex}

\begin{ex}%[De-chuan-hoa-so-3]%[Nguyễn Hoài Nam]%[0D0V2-4]
	Lớp $11A6$ có $45$ học sinh gồm $20$ nam và $25$ nữ. Chọn ngẫu nhiên $8$ bạn đi học đối tượng Đoàn, Tính xác suất để $8$ bạn được chọn có cả nam và nữ (làm tròn kết quả đến chữ số thập phân thứ hai).
	\shortans[oly]{$0{,}99$}
	\loigiai{
		Ta có $n(\Omega) = \mathrm{C}_{45}^8$.\\
		Gọi $A$ là biến cố \lq\lq chọn ra $8$ bạn đi học đối tượng Đoàn sao cho trong đó có cả nam lẫn nữ \rq\rq. \\
		Suy ra $\overline{A}$ là biến cố: \lq\lq chọn ra $8$ bạn đi học đối tượng Đoàn sao cho trong đó chỉ có nam hoặc nữ\rq\rq.\\
		Chọn ra $8$ bạn học sinh bất kì trong số $45$ học sinh có $\mathrm{C}_{45}^8$ cách.\\
		Chọn ra $8$ bạn học sinh bất kì trong số $20$ học sinh nam có $\mathrm{C}_{20}^8$ cách.\\
		Chọn ra $8$ bạn học sinh bất kì trong số $25$ học sinh nữ có $\mathrm{C}_{25}^8$ cách.\\
		Vậy số cách chọn ra $8$ bạn đi học đối tượng Đoàn sao cho trong đó có cả nam lẫn nữ là
		$$\mathrm{C}_{45}^8-\mathrm{C}_{20}^8-\mathrm{C}_{25}^8=214\,345\,650\; \text{cách}.$$
		Xác suất là $P(A)  =\dfrac{n(A)}{n(\Omega)}= \dfrac{214\,345\,650}{\mathrm{C}_{45}^8} \approx 0{,}99$.
	}
\end{ex}

\begin{ex}%[De-chuan-hoa-so-3]%[Nguyễn Hoài Nam]%[0D0V2-5]
	Một hộp đựng $15$ viên bi, trong đó có $7$ viên bi xanh và $8$ viên bi đỏ. Lấy ngẫu nhiên ba viên bi (không kể thứ tự ra khỏi hộp). Tính xác suất để trong ba viên bi lấy ra có ít nhất một viên bi đỏ (làm tròn kết quả đến chữ số thập phân thứ hai).
	\shortans[oly]{$0{,}92$}
	\loigiai{
		Số phần tử của không gian mẫu $n(\Omega)=\mathrm{C}_{15}^3=455$.\\
		Gọi $A$ là biến cố \lq\lq trong ba viên bi lấy ra có ít nhất một viên bi đỏ \rq\rq.\\
		Suy ra $\overline{A}$ là biến cố \lq\lq trong ba viên bi lấy ra không có bi đỏ nào\rq\rq.\\
		Chọn $3$ bi xanh trong $7$ bi xanh có $\mathrm{C}_7^3=35$ cách.\\
		$$\Rightarrow n(\overline{A})=35\Rightarrow \mathrm{P}(\overline{A})=\dfrac{n(\overline{A})}{n(\Omega)}=\dfrac{1}{13}.$$
		Vậy $\mathrm{P}(A)=1-\mathrm{P}(\overline{A})=\dfrac{12}{13} \approx 0{,}92$.
	}
\end{ex}
\begin{ex}%[De-chuan-hoa-so-3]%[Nguyễn Hoài Nam]%[0H9H4-2]
	Viết phương trình đường tròn $(C)\colon x^2+y^2-2ax-2by+c=0$ biết $(C)$ đi qua hai điểm $A(3;3)$, $B(7;-1)$  và có tâm nằm trên đường thẳng $d\colon x+y-2=0$. Khi đó $a+b+c$ bằng bao nhiêu?
	\shortans[oly]{$-4$}
	\loigiai{
		Vì $(C)$ đi qua hai điểm $A(3;3)$, $B(7;-1)$  và có tâm nằm trên đường thẳng $d\colon x+y-2=0$ nên ta có hệ phương trình\\
		$\heva{&3^2+3^2-2a\cdot3-2b\cdot3+c=0\\&7^2+(-1)^2-2a\cdot7-2b\cdot(-1)+c=0\\&a+b-2=0}\Leftrightarrow\heva{&-6a-6b+c=-18\\&-14a+2b+c=-50\\&a+b=2}\Leftrightarrow\heva{&a=3\\&b=-1\\&c=-6.}$\\
		Vậy phương trình đường tròn cần tìm là
		$(C)\colon x^2+y^2-6x+2y-6=0.$
	}
\end{ex}
\begin{ex}%[De-chuan-hoa-so-3]%[Nguyễn Hoài Nam]%[0D0V2-6]
	Cho đa giác đều $20$ đỉnh $A_1.A_2\ldots A_{20}$ nội tiếp đường tròn tâm $O$. Chọn ngẫu nhiên ba đỉnh của đa giác đó, tính xác suất để ba đỉnh được chọn là ba đỉnh của một tam giác vuông (làm tròn kết quả đến chữ số thập phân thứ hai).
	\shortans[oly]{$0{,}16$}
	\loigiai{
		Gọi $A$ là biến cố \lq\lq ba đỉnh được chọn là ba đỉnh của một tam giác vuông\rq\rq.\\
		Số cách chọn ra một tam giác bất kì là số cách chọn ra $3$ đỉnh trong $20$ đỉnh của đa giác đã cho\\
		Suy ra $n(\Omega)=\mathrm{C}_{20}^3=1\,140$. \\
		Tam giác vuông nội tiếp đường tròn có một cạnh là đường kính.\\
		Chọn ra $2$ đỉnh là đường kính của đường tròn có $10$ cách.\\
		Chọn ra $1$ đỉnh trong $18$ đỉnh còn lại có $18$ cách.\\
		$\Rightarrow n(A)=10\cdot18=180$.\\
		Vậy $\mathrm{P}(A)=\dfrac{180}{1140}=\dfrac{3}{19} \approx 0{,}16$.
	}
\end{ex}

\Closesolutionfile{ans}
\indapan{7}{ans/ans-cau-KQ-1}